\section{Introduction}
\label{sec:introduction}

Multi-agent applications express a task as a directed acyclic graph of model
invocations. Each Agent can be called more than once, and later calls extend
a history of instructions, model outputs, tool results, and messages from
other Agents. A serving worker turns that history into a KV cache. When the
worker fails, the cache is gone. If no copy survives, the replacement replica
must prefill the retained history before it can emit the next token
\cite{strati2024dejavu,cao2026lumen}. Section~\ref{sec:motivation} measures
one extreme of that cost: at one million tokens, full reconstruction of
Qwen2.5-7B takes 571.2\,s, while reusing the same KV state from HBM, host
DRAM, or shared storage is ready in 1.3\,s, 2.7\,s, or 10.2\,s. The KV state
itself occupies 57.9\,GB. Protecting every active Agent at that scale is not
free.

The cost is also uneven. A delay on a branch that still has workflow slack
need not move the join, whereas the same delay on the critical branch does
\cite{lin2024parrot}. Native sparse attention concentrates historical accesses
on a subset of blocks \cite{jiang2024minference,huang2026nosav2}, which
suggests that a recovery copy might keep only those blocks in the fastest
memory. Finite peer memory then has to decide which copies stay. \system is
built from these three observations. A Tracker estimates, at each invocation
submission, how long remains until the Agent finishes, and turns that
estimate into guarded slack. KV Master stores a complete DRAM checkpoint for
Agents whose full reconstruction exceeds slack by more than a threshold, and
it additionally places the Top-$K$ blocks of the first $L$ layers in peer
HBM. An Adapter changes only the total HBM quota. The mechanisms are
specified in Section~\ref{sec:design}. They are not, by themselves, evidence
that the selection improves workflow completion.

Captured NOSA, NSA, and MiniCPM forwards arrive online on the recorded
BLSC A100-PCIe configuration. With no failure, \system finishes every
workflow at the same time as a restart. One GPU failure during a
24{,}576-token workflow shortens its retry, and every later workflow, by
1{,}585.4\,ms. Workflows released before the failure are unchanged. The HBM
blocks move completion by at most 0.217\,ms.

The contributions are the design and the evidence from this evaluation.
\begin{itemize}
\item A recovery design that keeps complete peer-DRAM checkpoints for
Agents whose estimated reconstruction exceeds slack. It also specifies a
Top-$K{\times}L$ peer-HBM tier and a quota.
\item An online evaluation on the recorded BLSC A100-PCIe configuration.
With no failure, completion time matches a restart. One GPU failure during
a 24{,}576-token workflow shortens that workflow and every later one by
1{,}585.4\,ms, and the saving grows with prompt length.
\item A result that the HBM blocks placed for the interrupted Agent move
completion by at most 0.217\,ms beyond the DRAM checkpoint, and that
$K{=}16$ covers 4--8\% of distinct blocks.
\end{itemize}
